\chapter{Implementing the Task API}
\chapterintro{The framework and repository meet in a thin application layer. Handlers translate HTTP inputs into validated repository calls and translate results into resource representations.}

\begin{objectives}
\begin{itemize}
  \item Build CRUD handlers without leaking HTTP into the repository.
  \item Validate create, replace, patch, path IDs, and list queries consistently.
  \item Run the API and exercise it with a standard-library client.
  \item Trace one create and one update through every layer.
\end{itemize}
\end{objectives}

\section{Application construction}

\code{build\_app(database\_path)} initializes the schema, creates a repository,
creates an Application, registers handlers, and returns it. The function makes
tests easy: every test supplies a temporary path and receives a fully configured
application without starting a socket.

The complete application layer follows. Read it once from top to bottom, then
return to the route-by-route discussion.

\codefile{../code/final/task_api/app.py}

\section{Create flow}

The POST handler requires a declared nonempty body, decodes JSON, validates a
full task document, supplies timestamps, and calls \code{repository.create}.
It returns the created dictionary as JSON with 201 and Location.

The server-generated fields never come from the request. This prevents a client
from choosing another task's identity or falsifying creation history. Unknown
input names fail validation, so a client typo such as \code{prioriy} is visible.

\section{Read and list flow}

The individual GET converts the captured ID. Invalid or nonpositive IDs receive
the same task-not-found category as absent records, keeping the public resource
contract simple.

The collection GET validates every query name, repetition count, enum, date,
integer range, and search length before repository access. The list result
combines items with total and effective pagination values. A caller can construct
the next offset without inferring defaults.

\section{PUT and PATCH}

Both update handlers validate content before touching storage and set a new UTC
update timestamp. PUT calls full-document validation, which requires title and
supplies defaults. PATCH calls partial validation, which rejects empty input and
includes only provided fields.

The repository return value distinguishes missing ID. A validation error is 422;
a valid modification targeting an absent task is 404. These outcomes must not be
collapsed into a generic 400.

\section{Date and timestamp choices}

Due date is a calendar date with no timezone. Validation parses it and requires
that serializing the parsed date produces exactly the original text, excluding
alternate compact forms accepted by some library versions.

Audit timestamps represent instants and use UTC with a trailing \code{Z}, to
second precision. A production system might use database-generated timestamps,
higher precision, or a test-injectable clock. Here the helper keeps one format
throughout the application.

\section{Run the API}

From \filepath{code/final}:

\begin{wiremessage}
python -m task_api --database tasks.db
\end{wiremessage}

The command binds to \code{127.0.0.1:8080}, initializes the database, starts 16
workers, and enables bounded keep-alive. Optional arguments select host, port,
database path, worker count, and keep-alive behavior.

In a second terminal:

\begin{wiremessage}
python client_example.py
\end{wiremessage}

The client uses \code{http.client}, creates a task, reads and prints the response,
then reuses the connection to list matching tasks. Because it is standard-library
code, the core demonstration does not depend on curl or a third-party SDK.

\codefile{../code/final/client_example.py}

\section{Trace a PATCH}

For \method{PATCH /tasks/1} with \code{\{"status":"done"\}}:

\begin{enumerate}
  \item HTTP parser frames the declared body and parses the target.
  \item Router captures \code{task\_id = "1"}.
  \item Handler checks body, media type, UTF-8, and JSON.
  \item Partial validator accepts only the status field and adds no defaults.
  \item Handler adds \code{updated\_at}.
  \item Repository generates assignments from its allow-list and binds values.
  \item SQLite updates inside a transaction and returns the complete row.
  \item Response factory encodes the representation; server frames and sends it.
\end{enumerate}

If any step fails, ownership is clear. A malformed request never reaches routing;
an invalid enum never reaches SQL; a missing row becomes a resource 404; an
unexpected database error crosses the generic 500 boundary.

\begin{tryit}
Create three tasks with different status, priority, and due dates. Query combined
filters and pagination. Then send an unknown query parameter, repeated limit,
invalid date, list instead of object, unsupported media type, and nonexistent ID.
Record status and problem type for each.
\end{tryit}

\summarytitle

The Task application is thin translation code between HTTP and a repository.
It validates each input category, owns resource semantics and status choices,
protects server-generated fields, and returns stable representations/problems.
The complete backend runs without third-party Python packages.

\reviewtitle
\begin{enumerate}
  \item Why does \code{build\_app} accept a database path instead of using a fixed global database?
  \item Which fields can a client write, and which remain server-owned?
  \item Why are alternate date spellings rejected even if Python can parse them?
  \item At which step does a missing row become an HTTP 404?
\end{enumerate}

\begin{exercise}[End-to-end extension]
Add a \code{GET /tasks/summary} static route returning counts by status. Write
the contract first, add a repository aggregate with parameter-free SQL, ensure
the static route beats \code{/tasks/\{task\_id\}}, and test response/problem cases.
\end{exercise}

\begin{exercise}[Optimistic concurrency design]
Without implementing it, design ETag and \code{If-Match} support for updates.
Define validator generation, success and precondition-failure statuses, database
atomicity, and behavior when the field is absent.
\end{exercise}
